\chapter{General Information and Key Concepts}
\thispagestyle{fancy}
\justifying

This chapter provides the theoretical framework for our study, exploring the critical landscape of emergency medical systems and the specific mechanisms of patient triage. To understand the context of the proposed digital solution, it is essential to first define the roles of institutional relief organizations in Burkina Faso and the internationally validated protocols used to manage the influx of patients into emergency departments and pre-hospital care.

\section{Emergency Medical Services and the SAMU}
In Burkina Faso, the Service d'Aide M\'edicale Urgente (SAMU) was established as a cornerstone of the public health infrastructure to provide specialized and free medical care for critical patients before they reach the hospital environment. The mission of the SAMU is twofold: first, to coordinate medical responses through a centralized regulation center; and second, to deploy medicalized emergency units known as SMUR (Service Mobile d'Urgence et de R\'eanimation). These units are equipped with advanced medical technology and staffed by specialized personnel capable of performing life-saving procedures in the field. Consequently, the effectiveness of the entire emergency chain depends heavily on the accuracy of the regulation center, where the first decisions about patient priority are made.

\section{The Strategic Importance of Triage}
Medical triage is more than a simple sorting procedure; it is a strategic management tool designed to optimize limited healthcare resources and ensure the highest possible survival rates during emergencies or mass-casualty incidents. Historically, the triage concept evolved from military medicine and has been adapted for modern hospital and pre-hospital environments. The fundamental principle revolves around the idea that patients should be treated not only based on their arrival time but, more importantly, on the severity of their condition and their need for immediate intervention. In a high-pressure environment such as the SAMU regulation center, where callers provide varying levels of information, a standardized triage protocol serves as a common language between dispatchers and field teams.

\section{The Emergency Severity Index (ESI)}
To address the need for a reliable and objective triage framework, the Emergency Severity Index (ESI) was developed as a five-level clinical triage tool. Unlike traditional three-level systems, the ESI provides a more granular stratification by considering both patient acuity and the expected consumption of medical resources. The ESI algorithm is structured around four main decision points, leading to a priority ranking from 1 (immediate life threat) to 5 (least urgent). For instance, patients in cardiac arrest or with severe respiratory distress are categorized as ESI Level 1, requiring immediate resuscitation. Conversely, patients who are stable and require only simple diagnostic tests are assigned to higher levels. This objectivity is essential for minimizing the risk of clinical oversight and for standardizing care across different medical teams.

\section{Digital Decision Support in Modern Healthcare}
The integration of digital technology into clinical workflows has revolutionized the way medical decisions are made. A Digital Decision Support System (DDSS) for triage acts as an intelligent intermediary that guides the operator through the ESI algorithm, ensuring that no critical diagnostic criteria are omitted. By digitizing a clinical protocol, we can significantly reduce the cognitive load on healthcare workers, mitigate the effects of fatigue or stress, and provide an auditable record of each decision. Modern systems, built on robust architectures such as FastAPI and React, offer the performance and security required to handle sensitive patient data while providing intuitive interfaces that fit naturally into the fast-paced environment of an emergency regulation center.

\clearpage
